\clearpage
\section{Stable assignment with preserved store identity}
\label{sec:copying-assignment}
Figure~19(c) (PDF p.~38) now has a concrete local assignment rule.
The matched atom contraction tree is arbitrary, the new value is an arbitrary
stable hypernet, and all other atom-reference inputs remain available.
The rule updates the selected store's term connection, weakens the consumed
instance and old contents, and returns unit. It neither allocates a fresh
store nor deletes the old heap. This matters for Proposition~12.2 because
its arbitrary callback belongs to the full language and may use assignment,
even though the displayed terms contain no assignment themselves.

The development proves exact preservation of store-resolution costs, ordered
boundary preservation, safety, renaming and recursive padding. Seven-step
local executions are checked, including a payload using the unchanged
Figure~16 literal translation. The local rule is not yet integrated into
arbitrary focused-context matching or the complete copying UAM. In
particular, no whole-program assignment execution or callback equivalence
is claimed here. The proposition's terms and translation remain unchanged.

\Needspace{12\baselineskip}
\subsection{Interfaces and finite assignment frames}

\termdef{copying-assignment-interface}
Let $C:\diamond^{h+1}\Rightarrow\diamond$ be the matched atom tree and
$G:\star\Rightarrow O$ the new-value graph. The fragment's input list is
one term input followed by $h$ atom inputs. Its output list is one term
output, for the old contents, followed by $O$, preserving the new value's
external interface. The assigning instance supplies input zero of $C$;
the other inputs of $C$ remain external. Connections from the value's atom
outputs back into that tree can be made by a surrounding graph; these
fragment interfaces expose those cuts rather than assuming the references
are distinct cells. Such arbitrary ambient matching remains a separate
integration obligation.

The frame contains a tree hole and a value hole with exactly the signatures
of $C$ and $G$. Its redex edge tags are assignment, instance, tree, store and
value. The reduct adds separate tags for discarding the instance, discarding
the old contents, and returning unit, while retaining instance, tree, store
and value tags.

\lean{Copying/AssignmentFrames}{9}{9}{Assignment.inputs}

\lean{Copying/AssignmentFrames}{10}{10}{Assignment.outputs}

\lean{Copying/AssignmentFrames}{11}{11}{Assignment.treeHole}

\lean{Copying/AssignmentFrames}{12}{12}{Assignment.valueHole}

\lean{Copying/AssignmentFrames}{14}{16}{Assignment.BeforeEdge}

\lean{Copying/AssignmentFrames}{17}{19}{Assignment.AfterEdge}

Input positions are embedded by their numerical indices. The other-reference
list is the increasing list of successor positions. Value outputs are the
successor-indexed output positions, leaving output zero for the old contents.
Both lists are duplicate-free by injectivity of the relevant port constructor
and successor. Their membership lemmas follow from finite-range mapping:
a vertex belongs precisely when it has the corresponding indexed form.

\lean{Copying/AssignmentFrames}{21}{22}{Assignment.inputPort}

\lean{Copying/AssignmentFrames}{23}{24}{Assignment.otherInputs}

\lean{Copying/AssignmentFrames}{25}{26}{Assignment.outputPort}

\lean{Copying/AssignmentFrames}{27}{28}{Assignment.valueOutputs}

\lean{Copying/AssignmentFrames}{30}{31}{Assignment.otherInputs_nodup}

\lean{Copying/AssignmentFrames}{32}{33}{Assignment.mem_otherInputs}

\lean{Copying/AssignmentFrames}{34}{35}{Assignment.valueOutputs_nodup}

\lean{Copying/AssignmentFrames}{36}{37}{Assignment.mem_valueOutputs}

The redex has four internal vertices. Two term vertices are the assignment
arguments: the first feeds the instance, the second feeds $G$. The instance
produces an atom vertex feeding the tree's first input; the other atom
vertex connects the tree root to the store. The store's term output is
external output zero. The value hole's outputs occupy the remaining
external positions. The label, source and target functions specify every
one of these incidences in order.

The finite-net constructor checks unique incidence, complete coverage,
duplicate-free incidence lists and signature typing. Internal coverage is
four-way finite case analysis. External coverage splits an index into zero
or successor. Uniqueness separates the five edge tags and uses disjointness
of zero and successor output positions. Mapping each boundary occurrence
to its type gives the required input/output signature lists.

\lean{Copying/AssignmentFrames}{39}{44}{Assignment.beforeLabel}

\lean{Copying/AssignmentFrames}{46}{51}{Assignment.beforeSource}

\lean{Copying/AssignmentFrames}{53}{58}{Assignment.beforeTarget}

\lean{Copying/AssignmentFrames}{60}{122}{Assignment.beforeFrame}

In the reduct, the tree and instance keep their atom wiring. A new term
weakening feeds the consumed instance. The store's output now feeds the
new value at the second term vertex. A separate term weakening feeds old
output zero. Unit consumes the external term input and has no outputs.
The same four internal vertices suffice, and all $O$ outputs keep their
positions. The seven-edge constructor proves the same finite typing and
incidence obligations, using the same zero/successor distinctions.

\lean{Copying/AssignmentFrames}{124}{131}{Assignment.afterLabel}

\lean{Copying/AssignmentFrames}{133}{140}{Assignment.afterSource}

\lean{Copying/AssignmentFrames}{142}{149}{Assignment.afterTarget}

\lean{Copying/AssignmentFrames}{151}{211}{Assignment.afterFrame}

\Needspace{12\baselineskip}
\subsection{The stable compute rule}

\termdef{copying-assignment-rule}
Relabelling the primitive frames into allowance $d$ produces two patterns.
The filling functions replace the tree hole by $C$ at that allowance and
the value hole by $G$. Every other frame edge is filled by its original
primitive corolla. Simultaneous expansion gives actual finite typed redex
and reduct hypernets; no holes remain from the pattern itself. This uses
$G$ literally, without simplifying distributors inside a translated value.

\lean{Copying/Assignment}{8}{9}{Assignment.beforePattern}

\lean{Copying/Assignment}{10}{11}{Assignment.afterPattern}

\lean{Copying/Assignment}{13}{23}{Assignment.beforeFill}

\lean{Copying/Assignment}{25}{36}{Assignment.afterFill}

\lean{Copying/Assignment}{38}{39}{Assignment.before}

\lean{Copying/Assignment}{40}{41}{Assignment.after}

The local rule requires the existing full stability predicate on $G$,
including its label, output-type and input-reachability conditions. No
stability obligation is assumed of the entire heap. Rule inversion
exposes exactly a matched tree, a stable value, and the equalities to these
two concrete graph constructions. The proof is elimination and construction
of the single assignment-rule constructor.

\lean{Copying/Assignment}{44}{46}{Assignment.Rule}

\lean{Copying/Assignment}{48}{52}{Assignment.rule_iff}

Designated edge injections locate the original instance, store and each
new-value edge on both sides. Their labels are preserved exactly, by
unfolding expansion at zero and successor allowance. In particular, the
store occurrence remains a store and each value edge keeps its original
operation or box label.

\lean{Copying/Assignment}{54}{55}{Assignment.beforeInstance}

\lean{Copying/Assignment}{56}{57}{Assignment.afterInstance}

\lean{Copying/Assignment}{58}{59}{Assignment.beforeStore}

\lean{Copying/Assignment}{60}{61}{Assignment.afterStore}

\lean{Copying/Assignment}{62}{63}{Assignment.beforeValue}

\lean{Copying/Assignment}{64}{65}{Assignment.afterValue}

\lean{Copying/Assignment}{67}{69}{Assignment.value_labels}

\lean{Copying/Assignment}{70}{72}{Assignment.store_labels}

\Needspace{12\baselineskip}
\subsection{Exact preservation of store resolution}

The tree-edge injections and vertex insertion maps locate the preserved
atom tree inside either rule graph. Every tree label is its original
primitive label lifted to the allowance. An adjacent pair of tree edges
therefore yields an atom-routing hop on either side: the tree's label
classification establishes a forwarding edge, and mapping the shared
vertex proves adjacency in the expanded net.

\lean{Copying/AssignmentResolution}{9}{10}{Assignment.beforeTree}

\lean{Copying/AssignmentResolution}{11}{12}{Assignment.afterTree}

\lean{Copying/AssignmentResolution}{13}{14}{Assignment.beforeTreeVertex}

\lean{Copying/AssignmentResolution}{15}{16}{Assignment.afterTreeVertex}

\lean{Copying/AssignmentResolution}{18}{20}{Assignment.tree_labels}

\lean{Copying/AssignmentResolution}{22}{31}{Assignment.tree_hop_before}

\lean{Copying/AssignmentResolution}{33}{42}{Assignment.tree_hop_after}

The instance enters input zero of the same tree on both sides. Given an
edge incident to that input, its adjacency with the instance follows from
the explicit frame attachment. At the other end, the tree's root output
is attached to the designated store input. These entry and exit lemmas
hold for arbitrary contraction-tree presentations, not only canonical ones.

\lean{Copying/AssignmentResolution}{45}{52}{Assignment.enters_before}

\lean{Copying/AssignmentResolution}{54}{61}{Assignment.enters_after}

\lean{Copying/AssignmentResolution}{63}{72}{Assignment.root_before}

\lean{Copying/AssignmentResolution}{74}{83}{Assignment.root_after}

\termdef{copying-assignment-resolution}
The contraction-tree reachability theorem supplies an input-zero edge and
a path of some length $m$ from that edge to the root. Map this path through
the two tree embeddings, prepend the instance-entry hop and append the
root-to-store hop. This gives resolution to the designated store on both
sides in exactly $m+2$ hops. The general uniqueness theorem for store traces
then proves equality of the successful resolution counts in both directions.
It also proves that any successful resolution from the instance must end
at the designated store, before or after assignment. Stability of the new
value is not needed for these routing lemmas; the rule itself still requires it.

\lean{Copying/AssignmentResolution}{86}{93}{Assignment.resolves_both}

\lean{Copying/AssignmentResolution}{96}{102}{Assignment.resolution_count_iff}

\lean{Copying/AssignmentResolution}{105}{108}{Assignment.store_unique_before}

\lean{Copying/AssignmentResolution}{111}{114}{Assignment.store_unique_after}

\Needspace{12\baselineskip}
\subsection{What changes and what remains attached}

The value-vertex insertion map places $G$ in the reduct. The store target
is exactly the image of $G$'s term input. Every source and target list of
$G$ is mapped entrywise by this insertion, and output $i$ of $G$ remains
external output $i+1$. These equations follow from the expansion definition
and the frame's ordered value-output list. They prove reconnection of the
actual value graph, rather than an abstract update of a separate heap map.

\lean{Copying/AssignmentPreservation}{9}{10}{Assignment.afterValueVertex}

\lean{Copying/AssignmentPreservation}{12}{13}{Assignment.updated_content}

\lean{Copying/AssignmentPreservation}{15}{16}{Assignment.value_sources}

\lean{Copying/AssignmentPreservation}{18}{19}{Assignment.value_targets}

\lean{Copying/AssignmentPreservation}{21}{27}{Assignment.value_output_order}

Separate edge selectors identify the two new weakenings. The first has
term-weakening label and feeds external old-content output zero. The other
has the same label and feeds precisely the consumed instance's input.
Neither weakening is erased. For each other reference input $i+1$, the
tree-vertex map before and after assignment is the same external input
position. Thus the update changes a term connection while preserving the
atom routing and other reference attachments.

\lean{Copying/AssignmentPreservation}{29}{30}{Assignment.discardOld}

\lean{Copying/AssignmentPreservation}{31}{32}{Assignment.discardInstance}

\lean{Copying/AssignmentPreservation}{34}{37}{Assignment.old_content_weakened}

\lean{Copying/AssignmentPreservation}{39}{42}{Assignment.instance_weakened}

\lean{Copying/AssignmentPreservation}{44}{53}{Assignment.other_references_preserved}

An arbitrary continuation may be attached to all fragment outputs, including
the old contents at output zero. Serial composition preserves every edge
label of that continuation. The already established successful atom trace
also survives this attachment, so store-trace uniqueness prevents the
updated instance from resolving to a different surrounding cell. This
argument imposes no acyclicity condition on the continuation. It is a
structural and routing preservation result, not yet a full contextual
execution theorem for assignment.

\lean{Copying/AssignmentPreservation}{57}{58}{Assignment.withContents}

\lean{Copying/AssignmentPreservation}{60}{62}{Assignment.old_contents_retained}

\lean{Copying/AssignmentPreservation}{65}{70}{Assignment.attached_store_unique}

\Needspace{12\baselineskip}
\subsection{Search, update and unit return}

The two argument coordinates are the first and second internal term
vertices. Search first enters the instance argument. The instance returns
by backtracking; the assignment head then resumes at the second, new-value
argument. Once that argument returns, backtracking through assignment
requests rewrite focus at the external term input. After the assignment
rewrite resets search there, the unit edge returns immediately. These five
individual interaction lemmas check the actual ordered argument positions.

\lean{Copying/AssignmentExecution}{8}{9}{Assignment.instanceArgument}

\lean{Copying/AssignmentExecution}{10}{11}{Assignment.valueArgument}

\lean{Copying/AssignmentExecution}{13}{19}{Assignment.search_instance}

\lean{Copying/AssignmentExecution}{21}{26}{Assignment.instance_returns}

\lean{Copying/AssignmentExecution}{28}{31}{Assignment.search_value}

\lean{Copying/AssignmentExecution}{33}{36}{Assignment.ready}

\lean{Copying/AssignmentExecution}{38}{44}{Assignment.unit_returns}

\termdef{copying-assignment-subsystem}
A state keeps a concrete graph and its focus coordinate together. The local
transition relation contains intrinsic search and the stable assignment
compute rule. Finality is backtracking at an external input. Such a state
cannot search further, by the established search theorem, and cannot assign,
because assignment requires rewrite mode. This proves final normality for
the subsystem machine. It does not assert global determinism of overlapping
matches or supply the remaining language rules.

\lean{Copying/AssignmentExecution}{46}{46}{Assignment.State}

\lean{Copying/AssignmentExecution}{47}{50}{Assignment.Step}

\lean{Copying/AssignmentExecution}{52}{61}{Assignment.final_normal}

\lean{Copying/AssignmentExecution}{63}{66}{Assignment.machine}

Any constant corolla returns at the new-value argument in one search step.
For such a value the complete local run has seven transitions: enter the
instance, return from it, move to the value, return from it, request rewrite,
perform assignment, and return unit. The run uses the proved stability of
the constant, so no unproved operational premise is hidden in this example.

\lean{Copying/AssignmentExecution}{69}{76}{Assignment.constant_returns}

\lean{Copying/AssignmentExecution}{79}{86}{Assignment.constant_run}

\Needspace{12\baselineskip}
\subsection{Renaming, padding and safety}

An ordered isomorphism of matched atom trees and a deep isomorphism of
value graphs induce deep isomorphisms on both rule sides. Simultaneous
expansion assembles identity isomorphisms on the primitive corollas, the
lifted tree isomorphism in the tree slot, and the supplied value isomorphism
in the value slot. Boundary order is unchanged.

\lean{Copying/AssignmentTransport}{8}{18}{Assignment.before_congr}

\lean{Copying/AssignmentTransport}{20}{32}{Assignment.after_congr}

Padding has an explicit graph-isomorphism proof for each side. The maps
keep frame vertices fixed and map each component edge or internal vertex
within its original component tag. Their inverses use the same tags. Labels
agree under the recursive padding map. Incidence follows by composing list
maps; the host-vertex lemma checks that padding commutes with attachment to
each frame vertex. Empty weakening sources and the empty unit target are
handled directly. The uniform relabelling presentation of padding is then
identified with the existing recursive raise operation, yielding deep
isomorphisms for the original assignment graphs.

\lean{Copying/AssignmentTransport}{34}{139}{Assignment.before_raise_graph}

\lean{Copying/AssignmentTransport}{141}{143}{Assignment.before_raise}

\lean{Copying/AssignmentTransport}{145}{261}{Assignment.after_raise_graph}

\lean{Copying/AssignmentTransport}{263}{265}{Assignment.after_raise}

Assignment up to representation isomorphism retains a witness tree and a
stable value, together with deep isomorphisms to the literal redex and
reduct. A raw rule gives this immediately. Raising transports the witnesses
using the preceding two graph laws and the proved preservation of stability.
Induction gives forward preservation for every finite padding amount.
No determinism hypothesis is used to infer padding invariance.

\lean{Copying/AssignmentTransport}{268}{269}{Assignment.UpToIso}

\lean{Copying/AssignmentTransport}{270}{272}{Assignment.rule_upToIso}

\lean{Copying/AssignmentTransport}{273}{277}{Assignment.upToIso_raise}

\lean{Copying/AssignmentTransport}{278}{282}{Assignment.upToIso_pad}

For any primitive-label predicate $P$, redex safety follows from safety of
$G$, classification of the tree's labels, and $P$ on assignment, instance
and store. Reduct safety additionally requires term weakening and unit.
The proof splits the finite frame components and applies simultaneous
expansion's label law. Excluding focus and hole labels gives both properties
on both sides whenever $G$ is stable. Conversely, the updated whole heap
is not stable: its designated store label violates the stable-label
condition. This prevents mistakenly treating the entire post-update graph
as the stable argument required by a later compute rule.

\lean{Copying/AssignmentSafety}{8}{21}{Assignment.before_labels}

\lean{Copying/AssignmentSafety}{23}{38}{Assignment.after_labels}

\lean{Copying/AssignmentSafety}{40}{51}{Assignment.stable_assignment_safe}

\lean{Copying/AssignmentSafety}{54}{60}{Assignment.after_not_stable}

\Needspace{12\baselineskip}
\subsection{Translated payloads and concrete update regressions}

The literal-edge selector identifies the natural-number edge in the original
closed Figure~16 constant translation, including its retained empty
distributors. That edge proves a one-step return at the value argument.
Substituting this return and the already proved stability of the translated
literal gives the same seven-step local assignment execution. Only the
payload is claimed to be this exact translated term; the whole local rule
fragment is not being presented as a completed translation of a source
assignment context.

\lean{Copying/AssignmentExamples}{8}{8}{Assignment.literalEdge}

\lean{Copying/AssignmentExamples}{10}{15}{Assignment.translated_value_returns}

\lean{Copying/AssignmentExamples}{18}{25}{Assignment.translated_value_run}

The updated-heap regression attaches the original translated old literal
to the reduct's old-content output, while the new payload is the original
translated new literal. Edge selectors locate the updated store, new
literal, old literal and old-content weakening. Label equations establish
which numerical value each edge represents. The store target equals the
new literal's source, and the weakening target equals the old literal's
source. The two literal occurrences are disjoint because they belong to
different serial components, even when their numerical values coincide.

\lean{Copying/AssignmentExamples}{28}{29}{Assignment.updatedHeap}

\lean{Copying/AssignmentExamples}{30}{31}{Assignment.updatedCell}

\lean{Copying/AssignmentExamples}{32}{33}{Assignment.updatedValue}

\lean{Copying/AssignmentExamples}{34}{35}{Assignment.discardedValue}

\lean{Copying/AssignmentExamples}{36}{37}{Assignment.oldWeakening}

\lean{Copying/AssignmentExamples}{39}{41}{Assignment.updated_literal}

\lean{Copying/AssignmentExamples}{43}{46}{Assignment.discarded_literal}

\lean{Copying/AssignmentExamples}{49}{53}{Assignment.old_new_disjoint}

Instantiating with the weakening-free binary atom tree tests a noncanonical
match with another exposed reference. Updating literal one to literal two
has exactly the store-to-two and weakening-to-one connections just proved.
This verifies concrete mutation while retaining the discarded graph for
later garbage-collection reasoning.

\lean{Copying/AssignmentExamples}{57}{62}{Assignment.one_to_two}

Validation checks 326 build jobs and the three axiom audits: 6 abstract,
477 graph and 576 copying results. Only standard Lean logical axioms occur.
Estimated completion of Proposition~12.2 is about 65\% by mathematical effort,
with broad uncertainty of 50--75\% (26 October 2026). Assignment construction
is removed from Appendix~B.1's pending compute rules. Atom equality,
arithmetic, application, full-machine integration and template robustness
remain, as do the final Figure~22 chain and semantic weakening argument.

